\section{Sessions --- save and resume}
\label{sec:session}

pyALDIC can serialize your entire setup (image folder reference,
parameters, Regions of Interest, physical units) into a single
JSON file so you can close the application and pick up exactly
where you left off.

\subsection{File format}

\begin{itemize}
  \item Extension: \texttt{.aldic.json}
  \item Single-file portable: the Regions of Interest are stored
        inline as base64-encoded PNG.
  \item Versioned: the top-level key \texttt{schema\_version}
        protects against silent misinterpretation of newer schemas.
\end{itemize}

\subsection{What is saved}

\begin{itemize}
  \item Image folder path and the ordered list of filenames.
  \item Every per-frame Region of Interest mask.
  \item Core DIC parameters: subset size, step, search range,
        tracking mode, reference-update policy, solver choice,
        refinement settings, initial-guess mode.
  \item Physical units (pixel size, unit, frame rate).
\end{itemize}

\subsection{What is not saved}

\begin{itemize}
  \item Pipeline results (re-run with Run --- cheap vs.\ re-
        configuring).
  \item Run state / progress / elapsed time.
  \item Transient visualization preferences (colormaps, opacity,
        colour ranges).
\end{itemize}

\subsection{Menu}

\begin{description}[leftmargin=1.2cm,style=nextline]
  \item[File $\to$ Save Session\ldots] (\texttt{Ctrl+S})
    Opens a file dialog. Suggested extension
    \texttt{*.aldic.json}. Overwrites if the file exists.
  \item[File $\to$ Open Session\ldots] (\texttt{Ctrl+O})
    Opens a file dialog. Filter: \texttt{pyALDIC Session
    (*.aldic.json)}.
\end{description}

\subsection{Opening a session whose images moved}

If the image folder path in the session no longer exists, pyALDIC
does \emph{not} fail the load. Instead it:
\begin{enumerate}
  \item Restores parameters and Regions of Interest.
  \item Logs a warning: ``Image folder ... no longer exists;
        parameters and Regions of Interest were restored but images
        must be re-selected manually.''
  \item Leaves the image list empty.
\end{enumerate}

You can then drop a new folder of the same images and the Regions
of Interest will still be aligned.

\begin{tipbox}
This behaviour is deliberate so you can share a \texttt{.aldic.json}
file with a colleague whose images live at a different path.
They re-drop the images; all parameters and Regions of Interest
come along.
\end{tipbox}

\subsection{Corrupt or future-version files}

Malformed JSON, non-object root, or an unknown
\texttt{schema\_version} raise a modal error dialog. The fatal-
error path is used so you cannot miss the failure. The application
keeps whatever state it had before the attempted load; nothing is
silently replaced.
